% Formulae retain the paper's established 1.1.question.formula indices.
\numberwithin{equation}{enumii}
\renewcommand{\theenumii}{\thesubsection.\arabic{enumii}}
\setcounter{enumii}{1}\setcounter{equation}{0}

\paragraph*{Quadratic equations}
For $ax^2+bx+c=0$, where $a\ne0$, the discriminant and roots are

\par\noindent\textbf{1.1.1}\quad
\begin{equation}
\Delta=b^2-4ac,\qquad x=\frac{-b\pm\sqrt{\Delta}}{2a}.
\label{eq:quad-general}
\end{equation}
Equal real roots occur when $\Delta=0$. Always check that the leading coefficient is nonzero before treating an expression as a quadratic equation.


\setcounter{enumii}{2}\setcounter{equation}{0}
\paragraph*{Right-triangle trigonometry and motion}
For an angle $\theta$ in a right triangle,

\par\noindent\textbf{1.1.2}\quad
\begin{equation}
\tan\theta=\frac{\text{opposite}}{\text{adjacent}},\quad
\sin\theta=\frac{\text{opposite}}{\text{hypotenuse}},\quad
\cos\theta=\frac{\text{adjacent}}{\text{hypotenuse}}.
\label{eq:trig-ratios}
\end{equation}
Useful values are $\tan30^\circ=1/\sqrt3$, $\tan60^\circ=\sqrt3$, $\cot30^\circ=\sqrt3$, $\cot60^\circ=1/\sqrt3$, $\sin30^\circ=1/2$, and $\cos60^\circ=1/2$. For a right triangle,
\[c^2=a^2+b^2,\qquad \|(x,y)^T\|=\sqrt{x^2+y^2}.\]

\setcounter{enumii}{3}\setcounter{equation}{0}
\paragraph*{Probability}
For equally likely outcomes,

\par\noindent\textbf{1.1.3}\quad
\begin{equation}
P(E)=\frac{n(E)}{n(S)}.
\label{eq:prob-classical}
\end{equation}
When two different dice are tossed, there are $6\cdot6=36$ equally likely ordered outcomes; exchanging the dice generally gives a different outcome.

\paragraph*{Circles and tangents}
A radius is perpendicular to the tangent at the point of contact. Tangents drawn from the same external point have equal lengths. The angle between a tangent and a chord equals the angle in the alternate segment; the angle at the centre is twice the angle at the circumference standing on the same arc.
\setcounter{enumii}{4}\setcounter{equation}{0}
The tangent--chord theorem is
\par\noindent\textbf{1.1.4}\quad
\begin{equation}
\theta_{\mathrm{tangent-chord}}=\theta_{\mathrm{alternate\ segment}}.\tag{1.1.4.1}\label{eq:tangent-chord}
\end{equation}
\setcounter{enumii}{6}\setcounter{equation}{0}
For a triangle with inradius $r$ and semiperimeter $s$,
\par\noindent\textbf{1.1.6}\quad
\begin{equation}
[ABC]=rs,\qquad s=\frac{AB+BC+CA}{2}.
\label{eq:triangle-inradius}
\end{equation}
Tangents from the same external point have equal lengths:
\begin{equation}
PA=PB.\tag{1.1.6.2}\label{eq:equal-tangents}
\end{equation}

\setcounter{enumii}{8}\setcounter{equation}{0}
\paragraph*{Arithmetic progressions}
For first term $a$ and common difference $d$,

\par\noindent\textbf{1.1.8}\quad
\begin{equation}
t_n=a+(n-1)d,\qquad
S_n=\frac n2\bigl(2a+(n-1)d\bigr)=\frac n2(a+t_n).
\label{eq:ap}
\end{equation}
If the unknown parameters satisfy two linear equations, write them as an augmented matrix and use Gaussian elimination (elementary row operations) to solve the system.

\paragraph*{Coordinate geometry as linear algebra}
Two nonzero vectors are perpendicular exactly when their inner product is zero.
For a right angle at $B$,
\par\noindent\textbf{1.1.9}\quad
\begin{equation}
\overrightarrow{BA}\cdot\overrightarrow{BC}=0.\tag{1.1.9.1}\label{eq:right-angle-dot}
\end{equation}

\setcounter{enumii}{10}\setcounter{equation}{0}
For $A(x_1,y_1)$, $B(x_2,y_2)$, and $C(x_3,y_3)$, collinearity is equivalent to

\par\noindent\textbf{1.1.10}\quad
\begin{equation}
\mydet{x_2-x_1&y_2-y_1\\x_3-x_1&y_3-y_1}=0.
\label{eq:collinear-det}
\end{equation}

\paragraph*{Arithmetic progressions}
\setcounter{enumii}{11}\setcounter{equation}{0}
For an arithmetic progression,
\par\noindent\textbf{1.1.11}\quad
\begin{equation}
 t_n=a+(n-1)d,\qquad S_n=\frac n2\left(2a+(n-1)d\right).\tag{1.1.11.1}\label{eq:ap-question-11}
\end{equation}

Uniform motion obeys

\setcounter{enumii}{13}\setcounter{equation}{0}
\par\noindent\textbf{1.1.13}\quad
\begin{equation}
\text{distance}=\text{speed}\times\text{time},\qquad
\text{average speed}=\frac{\text{total distance}}{\text{total time}}.
\label{eq:speed-time}
\end{equation}

\setcounter{enumii}{14}\setcounter{equation}{0}
\paragraph*{Coordinate geometry as linear algebra}
Represent $P(x,y)$ by $\vec p=\myvec{x\\y}$. Then

\par\noindent\textbf{1.1.14}\quad
\begin{equation}
\overrightarrow{PQ}=\vec q-\vec p,\qquad
\langle\vec u,\vec v\rangle=u_xv_x+u_yv_y.
\label{eq:vector-basics}
\end{equation}
\setcounter{enumii}{14}\setcounter{equation}{1}
If $P$ divides $AB$ internally in the ratio $AP:PB=m:n$, then



\begin{equation}
\vec p=\frac{n\vec a+m\vec b}{m+n}=\vec a+\frac{m}{m+n}(\vec b-\vec a).
\label{eq:section-vector}
\end{equation}

\paragraph*{Probability}
\setcounter{enumii}{15}\setcounter{equation}{0}
If $E_1,E_2,E_3$ form an exhaustive set of mutually exclusive events, then
\par\noindent\textbf{1.1.15}\quad
\begin{equation}
P(E_1)+P(E_2)+P(E_3)=1.\tag{1.1.15.1}\label{eq:three-events}
\end{equation}


\setcounter{enumii}{16}\setcounter{equation}{0}
\paragraph*{Areas and solids}
For a sector of angle $\theta$ degrees and radius $r$,

\par\noindent\textbf{1.1.16}\quad
\begin{equation}
A_{\mathrm{sector}}=\frac{\theta}{360^\circ}\pi r^2,\qquad
A_{\mathrm{segment}}=A_{\mathrm{sector}}-A_{\mathrm{triangle}}.
\label{eq:sector-segment}
\end{equation}

\setcounter{enumii}{17}\setcounter{equation}{0}
The curved surface areas of a cylinder and a cone are

\par\noindent\textbf{1.1.17}\quad
\begin{equation}
A_{\mathrm{cyl}}=2\pi rh,\qquad A_{\mathrm{cone}}=\pi r\ell,\qquad \ell=\sqrt{r^2+h^2}.
\label{eq:cone-formula}
\end{equation}

\paragraph*{Arithmetic progressions}
\setcounter{enumii}{22}\setcounter{equation}{0}
The sum of terms $m$ through $n$ is
\par\noindent\textbf{1.1.22}\quad
\begin{equation}
S_{m\ldots n}=S_n-S_{m-1}.\tag{1.1.22.1}\label{eq:ap-subsum}
\end{equation}

\paragraph*{Right-triangle trigonometry and motion}
\setcounter{enumii}{23}\setcounter{equation}{0}
For two parts of a journey, total time is
\par\noindent\textbf{1.1.23}\quad
\begin{equation}
t_1+t_2=\frac{d_1}{v_1}+\frac{d_2}{v_2}.\tag{1.1.23.1}\label{eq:two-part-time}
\end{equation}

\setcounter{enumii}{27}\setcounter{equation}{0}
For an object at fixed height $h$ observed from level ground, its horizontal range is $x=h\cot\theta$.
In this situation,
\par\noindent\textbf{1.1.27}\quad
\begin{equation}
x=h\cot\theta.\tag{1.1.27.1}\label{eq:horizontal-range}
\end{equation}

\paragraph*{Probability}
\setcounter{enumii}{28}\setcounter{equation}{0}
For events $A$ and $B$,

\par\noindent\textbf{1.1.28}\quad
\begin{equation}
P(A\cup B)=P(A)+P(B)-P(A\cap B),\qquad P(A^c)=1-P(A).
\label{eq:prob-union}
\end{equation}

\paragraph*{Coordinate geometry as linear algebra}
\setcounter{enumii}{29}\setcounter{equation}{0}
The triangle area is half the absolute determinant:

\par\noindent\textbf{1.1.29}\quad
\begin{equation}
[ABC]=\frac12\left|\mydet{x_2-x_1&y_2-y_1\\x_3-x_1&y_3-y_1}\right|.
\label{eq:triangle-det}
\end{equation}

\paragraph*{Areas and solids}
\setcounter{enumii}{30}\setcounter{equation}{0}
A shaded region between an outer and an inner shape has area
\par\noindent\textbf{1.1.30}\quad
\begin{equation}
A_{\mathrm{shaded}}=A_{\mathrm{outer}}-A_{\mathrm{inner}}.\tag{1.1.30.1}\label{eq:shaded-area}
\end{equation}

\setcounter{enumii}{31}\setcounter{equation}{0}
The corresponding volume formula for a cylinder is

\par\noindent\textbf{1.1.31}\quad
\begin{equation}
V_{\mathrm{cyl}}=\pi r^2h.
\label{eq:cylinder-formula}
\end{equation}

\setcounter{enumii}{31}\setcounter{equation}{1}
For a hemisphere and sphere,


\begin{equation}
V_{\mathrm{hemisphere}}=\frac23\pi r^3,\qquad V_{\mathrm{sphere}}=\frac43\pi r^3,
\qquad A_{\mathrm{sphere}}=4\pi r^2.
\label{eq:sphere-formula}
\end{equation}

\setcounter{enumii}{31}\setcounter{equation}{2}
For a cylinder with two hemispherical cavities of radius $r$ scooped from its ends,



\begin{equation}
V_{\mathrm{remaining}}=\pi R^2h-\frac43\pi r^3.
\label{eq:scooped-cylinder}
\end{equation}

